\documentclass[10pt,a4paper]{amsart}
\usepackage[margin=19mm]{geometry}
\usepackage{amsmath,amssymb,mathtools}
\usepackage[scr=rsfs,cal=euler]{mathalfa}
\usepackage{xfrac}
\usepackage[unicode,hidelinks,bookmarks=false]{hyperref}
\newcommand{\ie}{i.e.\ }
\newcommand{\cf}{cf.\ }
\newcommand{\resp}{respectively\ }
\newcommand{\Subsection}{\subsection}
\providecommand{\nobreakdash}{\nobreak\hbox{-}}
\setlength{\emergencystretch}{2em}
\multlinegap0pt
\usepackage{bm}
\usepackage{bbm}
\usepackage{dsfont}
\usepackage{xspace}
\usepackage{cancel}
\usepackage{mathtools}
\usepackage{amsthm}
\makeatletter
\@ifundefined{thm}{\newtheorem{thm}{Thm}}{}
\@ifundefined{cor}{\newtheorem{cor}[thm]{Corollary}}{}
\@ifundefined{lem}{\newtheorem{lem}[thm]{Lemma}}{}
\makeatother
\DeclareMathOperator{\dB}{dB}
\newcommand{\dbf}{{f\!\!f}}
\newcommand{\dbt}{{t\!t}}
\title[d-Boolean algebras and their bitopological representation]{d-Boolean algebras and their bitopological representation}
\author{Hang Yang, Dexue Zhang}
\date{Source: arXiv:2505.17806v4 (CC BY 4.0)}
\begin{document}
\maketitle

\noindent\emph{Source and licence.} d-Boolean algebras and their bitopological representation. Version arXiv:2505.17806v4 (CC BY 4.0), distributed under the Creative Commons Attribution 4.0 International licence, as recorded in the arXiv licence metadata for this exact version and shown on the abstract page https://arxiv.org/abs/2505.17806v4. Full rights notice: licenses/arxiv-2505.17806-CC-BY-4.0.txt.\par\smallskip

\noindent\textit{Editorial scope.} Counted unit: the Lem whose source label is \texttt{complete dBA vs extremal discon} in the article (source0.tex lines 497--499, source label \texttt{complete dBA vs extremal discon}), delivered together with the article's own complete original proof of it (source0.tex lines 501--516, one \texttt{proof} environment of 16 lines). No mathematics is written, repaired or re-derived: statement and proof are the article's own text line for line. Required context retained uncounted: dBA via order reversing (lines 364--364). Every definition, display and label of those lines is retained unchanged. Omitted from this package and not redistributed: 1--363, 365--496, 500--500, 517--813. Nothing inside a retained excerpt is omitted or altered: every retained body is delivered exactly as the article's lines are written, so no omission receipt of any kind accompanies this package. Citations: the retained excerpts carry no citation command at all, so no bibliography entry of the article is transcribed and no reference is redistributed. Reference adaptation: none was needed and none was made; every reference inside the delivered unit resolves inside it, so no unresolved reference or citation is produced. Printed slips kept literal: no printed word, symbol, hypothesis or number of the counted statement or of its proof was altered. Layout and class: the article's own class is \texttt{article} with packages including amsfonts, amssymb, amsmath, amscd, amsthm, mathrsfs, mathtools, bbm, latexsym, color, verbatim, enumitem, hyperref, xy; the wrapper is the lane's pinned \texttt{amsart} editorial wrapper at 10pt on A4, which restates the theorem-like environments the retained text needs and adds the article's own macro definitions that the retained text needs (\texttt{\textbackslash dB}, \texttt{\textbackslash dbf}, \texttt{\textbackslash dbt}), with extra wrapper packages (bm, bbm, dsfont, xspace, cancel, mathtools). The delivered numbering is the unit's own. Assets: no retained span contains a figure environment, a graphics-inclusion command, a PDF-page inclusion, a table or a TikZ picture; no image file is copied into this package, the package declares no asset dependency and no image file is redistributed with it. Licence: version arXiv:2505.17806v4 of this article is distributed under the Creative Commons Attribution 4.0 International licence, as recorded in the arXiv licence metadata for this exact version and shown on the abstract page https://arxiv.org/abs/2505.17806v4. A full rights notice is delivered at licenses/arxiv-2505.17806-CC-BY-4.0.txt. Characters: every retained excerpt is pure ASCII, so the delivered engine, XeLaTeX with its default Latin Modern text font, reports no missing character.
\begin{cor}\label{dBA via order reversing} The consistency predicate {\bf con} of each d-Boolean algebra is a Scott closed set under the information order. \end{cor} 

\begin{lem} \label{complete dBA vs extremal discon} Let $\mathcal{L}=(L;\dbt,\dbf;{\bf con},{\bf tot})$ be  d-frame.   If  $\mathcal{L}$  is extremally disconnected, then the d-Boolean algebra $\dB\mathcal{L}$ is complete.
The converse implication also holds provided that   $\mathcal{L}$ is d-zero-dimensional.  
\end{lem}

\begin{proof} Suppose $\mathcal{L}$ is extremally disconnected. To see that the d-Boolean algebra $\dB\mathcal{L}$ (i.e., the d-Boolean coreflection of $\mathcal{L}$) is complete, it suffices to show that both $B_+$ and $B_-$ are complete lattices. Since taking d-pseudo-complement defines an anti-tone Galois connection between the complete lattices $L_+$ and $L_-$, the subset $\{a^*:a\in L_+\}$  is a complete lattice under the information order. By extremal disconnectedness of $\mathcal{L}$, $B_-$ is equal to the set $\{a^*:a\in L_+\}$, so $B_-$ is a complete lattice. Likewise, $B_+$ is a complete lattice. 

Now we show that if $\mathcal{L}$ is d-zero-dimensional and the d-Boolean algebra $\dB\mathcal{L}$ is complete, then $\mathcal{L}$ is extremally disconnected. Here we only check that the d-pseudo-complement  of each element of $L_+$ is d-complemented. The case for elements of $L_-$ is similar. 

Let $a\in L_+$. Since every element of $B_-$ is d-complemented, it suffices to show that the d-pseudo-complement $a^*$ of $a$ belongs to $B_-$. To this end, we show that  $$a^*={\sup}_{B_-} \{b\in B_{-}:b\sqsubseteq a^*\},$$ where the symbol ${\sup}_{B_-}$ denotes  join  in the complete lattice $B_-$. We proceed with two steps. 

{\bf Step 1}.   $a^*\sqsubseteq{\sup}_{B_-} \{b\in B_{-}:b\sqsubseteq a^*\}$. 

Since $\mathcal{L}$ is d-zero-dimensional, $a^*$ is a join (in $L$) of d-complemented elements (which are necessarily in $B_-$), then $a^*=\sup_L \{b\in B_{-}:b\sqsubseteq a^*\}$, hence $a^*\sqsubseteq{\sup}_{B_-} \{b\in B_{-}:b\sqsubseteq a^*\}.$   

{\bf Step 2}.   $a^*\sqsupseteq{\sup}_{B_-} \{b\in B_{-}:b\sqsubseteq a^*\}$. 

It suffices to show that  $a\sqcup \big({\sup}_{B_-}\{b\in B_-:b\sqsubseteq a^*\}\big)\in{\bf con}.$ 
For each $x\in B_{+}$ with $x\sqsubseteq a$ and each $b\in B_{-}$ with $b\sqsubseteq a^*$, since $a\sqcup a^*\in{\bf con}$, then $x\sqcup b\in{\bf con}$, hence $x\sqcup b\in{\bf con}_{\dB}$. Since  $\dB\mathcal{L}$ is a d-Boolean algebra, the consistency predicate ${\bf con}_{\dB}$ is Scott closed in $(\dB L,\sqsubseteq)$ by Corollary \ref{dBA via order reversing}.  Since $\{b\in B_-:b\sqsubseteq a^*\}$ is directed under the information order $\sqsubseteq$, then $$x\sqcup\big({\sup}_{B_-}\{b\in B_-:b\sqsubseteq a^*\}\big)\in {\bf con}_{\dB}\subseteq{\bf con}$$ for all $x\in B_{+}$ with $x\sqsubseteq a$.   

Since $\mathcal{L}$ is d-zero-dimensional,   $a$ is equal to the join of $\{x\in B_{+}:x\sqsubseteq a\}.$   Since the set $\{x\in B_+:x\sqsubseteq a\}$ is  directed  and ${\bf con}$ is a Scott closed set of $(L,\sqsubseteq)$,   it follows that $$a\sqcup\big({\sup}_{B_-}\{b\in B_-:b\sqsubseteq a^*\}\big)\in{\bf con}, $$ as desired.  \end{proof}

\end{document}
